% !TeX program = xelatex
% !TEX encoding = UTF-8
\documentclass[11pt,a4paper]{article}
\usepackage{zhpl}
\usepackage{pgfplotstable}
\usepackage{calc}
\newtheorem{corollary}{Следствие}
\newtheorem{lemma}{Лемма}
\renewcommand{\floatpagefraction}{0.7}

% ---------- Данные для рисунков ----------
% Модельная (синтетическая) реализация, посчитанная заранее: 2600 участков для
% рис. 1 и 60000 участков для гистограммы рис. 4. Параметры модели — в разделах 2–5:
% ядро (t0, x0) = (0,40; 0,40), sigma_t = 0,07, sigma_x = 0,06; доля захваченных
% f = 0,30; kappa ~ U(0, kappa*); lambda ~ Beta(1,4; 1,0); доля захваченных в
% собственных состояниях явки eta = 0,45; веса уровней g_n = 1 : 4 : 8;
% число избирателей на участке N ~ U{800, ..., 2499}; зерно ГСЧ 20111204.
\pgfplotstableread{
t x
42.2 25.9
31.5 48.2
18.1 34.8
41.3 34.9
35.1 29.6
42.0 39.0
37.2 42.7
28.0 24.5
45.2 38.8
40.2 35.1
31.6 38.0
36.1 35.0
40.3 48.8
49.5 30.9
32.7 55.3
45.7 38.1
40.3 27.7
43.0 46.1
38.6 46.7
52.0 36.2
39.0 44.3
38.3 38.9
38.5 31.8
38.6 39.4
45.0 50.1
47.1 45.2
35.9 46.1
58.8 42.5
36.2 43.2
50.0 47.5
29.8 41.0
29.9 35.2
27.0 40.1
39.5 28.9
32.6 39.7
50.3 37.3
45.6 51.1
45.0 37.8
45.7 30.8
44.4 35.4
36.3 41.9
38.2 35.9
47.7 38.0
41.7 43.3
33.2 35.6
45.0 49.6
35.9 46.0
32.7 34.2
30.8 41.3
44.8 39.2
42.9 41.1
39.5 29.4
51.5 42.3
53.3 30.4
41.9 42.9
53.1 37.1
39.2 41.1
35.3 41.5
37.5 50.5
35.9 44.0
54.7 31.0
31.1 30.2
42.5 32.1
50.9 35.9
38.4 47.4
33.6 38.6
49.6 50.8
33.6 37.3
30.1 39.6
45.8 44.9
57.6 48.0
44.5 33.3
36.4 40.2
37.1 38.3
45.0 42.3
41.9 33.3
45.0 50.6
43.3 42.2
37.9 44.6
35.0 44.4
45.0 35.9
43.4 46.0
46.4 41.1
36.5 29.8
42.0 43.0
35.4 34.5
40.1 38.7
32.5 43.9
49.8 40.2
39.9 42.8
36.1 40.0
30.6 36.3
39.4 31.9
39.3 37.6
35.8 43.1
42.1 47.8
46.9 40.7
42.4 39.0
31.1 38.9
35.2 36.6
34.8 37.2
41.7 36.4
48.1 36.1
26.3 36.8
46.7 47.7
47.9 31.1
37.2 41.0
42.0 32.9
46.8 44.9
34.4 45.3
52.4 49.0
38.0 47.0
35.1 41.1
34.9 37.4
43.8 34.1
21.9 27.4
39.0 40.4
50.2 39.5
30.7 30.1
50.9 29.4
36.0 43.7
41.5 42.0
30.8 26.7
34.4 47.2
26.9 44.3
36.1 29.4
30.9 44.0
43.5 32.5
40.4 30.0
50.5 40.2
42.7 41.0
54.7 38.4
37.7 35.6
34.2 35.9
44.9 43.0
37.1 38.1
49.9 39.1
45.4 37.6
49.7 49.6
39.9 39.4
34.7 36.0
29.0 46.9
38.7 39.0
38.7 56.0
42.3 43.2
42.5 44.8
34.9 43.8
38.1 37.2
50.4 38.9
41.0 51.3
39.3 42.9
49.0 31.6
34.3 40.6
37.2 38.3
40.9 36.8
31.8 40.6
49.2 45.5
40.6 39.1
38.8 42.1
33.6 41.8
40.2 31.2
40.8 35.4
31.2 48.7
50.1 40.6
40.2 39.0
35.9 35.2
43.0 40.9
37.4 33.0
43.9 35.6
27.6 38.2
24.2 40.9
42.4 40.6
40.1 58.4
41.1 34.4
24.2 43.3
44.8 39.4
43.6 47.5
41.5 47.6
56.5 43.1
51.5 46.9
35.6 39.7
35.4 42.5
53.7 34.1
42.1 31.6
42.7 41.7
33.2 49.2
31.2 48.6
22.9 42.1
22.6 46.3
34.2 51.6
41.3 43.2
33.6 40.8
45.3 48.0
33.6 41.4
39.8 36.4
36.2 37.2
38.2 34.1
46.0 35.7
36.9 36.3
37.3 47.9
47.6 29.8
43.5 31.4
39.3 40.7
33.5 39.7
38.1 46.6
41.2 25.3
39.8 29.9
42.2 38.6
40.0 37.7
45.2 40.3
43.1 35.4
41.1 32.7
51.6 41.8
36.3 31.7
44.2 37.5
48.5 48.0
35.9 33.3
36.0 43.2
50.0 45.0
48.4 38.6
40.4 35.4
40.3 29.5
54.5 42.2
37.4 30.9
43.2 44.7
31.7 40.9
48.0 46.3
41.5 34.4
45.3 35.1
53.2 23.9
41.9 44.1
41.9 47.0
43.4 43.8
38.7 46.1
40.3 40.1
33.9 44.7
44.4 41.9
47.0 43.6
34.4 40.0
47.6 42.9
36.9 30.4
48.5 40.4
33.0 31.0
33.8 43.2
39.5 42.6
47.5 41.1
48.2 36.6
34.4 35.5
42.5 34.2
48.2 36.9
42.3 26.0
37.7 48.0
41.1 29.7
48.7 42.1
34.2 42.1
36.0 31.6
40.5 43.0
50.3 46.6
40.2 23.8
34.3 41.3
43.8 45.7
31.6 40.3
33.1 46.6
35.6 40.0
52.3 42.4
54.5 37.0
35.8 43.7
35.2 30.5
51.3 37.1
39.1 47.6
47.5 37.2
39.3 43.2
42.1 52.4
47.9 37.7
32.4 34.2
43.0 41.4
41.3 38.5
50.9 37.9
34.9 43.8
48.2 40.0
43.2 45.1
36.1 44.6
36.7 36.2
34.1 42.1
37.3 45.5
57.3 40.0
43.6 42.3
47.9 28.2
40.1 32.1
44.7 38.5
35.1 39.0
39.8 39.4
34.8 41.9
36.5 32.3
45.7 50.0
37.6 39.8
49.2 46.5
47.1 26.9
37.4 24.9
31.6 44.0
36.0 34.4
42.4 40.3
34.3 35.3
25.4 27.1
29.4 37.4
45.5 49.4
44.3 44.0
34.7 41.9
33.6 31.7
36.0 36.4
36.0 37.8
43.5 31.1
37.9 40.8
33.9 39.7
62.2 46.2
49.4 36.2
42.2 43.0
39.4 41.8
44.0 47.1
42.4 34.7
47.0 45.0
43.0 36.5
37.7 36.9
43.3 38.6
35.3 33.1
48.5 33.6
40.2 51.4
26.4 46.6
47.9 29.6
50.0 46.2
29.7 36.7
38.5 38.8
43.0 35.1
25.5 35.5
37.1 50.2
39.6 38.3
46.5 35.9
30.7 36.8
39.6 37.4
47.0 33.0
48.8 41.3
37.7 42.1
40.0 36.7
37.8 38.2
26.3 33.0
35.3 45.9
31.3 42.6
35.3 30.4
42.8 44.7
31.5 41.5
43.4 40.5
55.5 45.0
44.9 42.1
44.3 40.5
50.7 35.0
42.0 41.9
50.5 40.2
51.7 49.1
30.5 32.4
42.0 35.3
60.9 31.2
36.8 42.1
54.5 47.0
40.5 44.9
40.6 36.9
47.0 42.5
36.1 43.3
42.2 30.0
32.8 38.8
33.8 46.7
43.3 40.1
41.4 43.5
40.5 48.2
45.9 43.8
43.7 40.0
47.6 43.7
26.9 48.3
29.5 38.9
33.2 37.9
32.0 41.3
25.9 45.7
27.8 45.3
44.9 35.8
39.8 35.8
34.1 45.0
38.1 46.5
38.9 34.0
51.9 44.6
41.2 46.4
37.2 42.3
37.5 41.9
40.1 41.7
43.9 37.9
48.2 45.7
43.7 43.8
30.4 30.6
39.7 39.5
40.7 38.8
39.3 37.3
34.6 47.1
50.0 36.5
18.3 36.4
45.3 39.7
37.7 47.0
40.0 39.8
35.1 41.4
43.5 36.4
37.8 23.8
49.8 39.1
37.1 43.6
48.9 30.8
37.8 47.6
41.6 39.6
40.3 33.1
36.6 37.6
33.0 44.0
37.6 47.8
27.1 51.1
41.2 42.6
38.2 42.4
34.8 35.9
50.0 31.8
47.3 44.3
38.2 34.9
35.9 29.1
37.2 33.2
45.6 52.7
38.4 47.7
29.0 48.8
37.6 42.4
46.2 38.2
44.3 29.1
33.1 42.5
47.6 38.6
35.6 42.9
45.8 34.7
30.5 39.1
40.2 33.5
40.3 39.4
37.0 38.4
24.9 41.5
31.1 47.0
45.5 40.8
38.1 37.7
40.8 43.7
33.5 45.6
46.5 40.8
52.9 43.6
31.5 41.0
46.3 38.0
32.5 37.7
35.0 47.2
36.9 34.9
36.9 37.8
30.8 38.4
31.7 43.1
44.8 36.4
45.3 39.6
32.2 46.9
38.2 43.3
42.9 44.8
37.5 40.0
37.8 48.8
51.9 34.8
51.2 35.8
42.1 49.0
40.6 49.4
28.2 38.0
28.9 44.5
37.7 42.1
42.4 30.4
42.8 38.6
37.2 30.6
46.1 29.6
33.5 47.7
52.6 30.9
52.5 48.2
32.6 53.4
42.1 44.5
28.5 36.3
37.0 35.4
51.4 43.5
38.1 28.7
42.2 37.9
37.2 49.0
35.0 34.6
36.8 28.7
41.9 33.5
48.9 38.6
48.0 40.5
44.7 34.7
49.9 37.1
38.9 34.5
39.0 39.9
34.9 31.6
26.3 42.2
31.9 41.1
35.7 34.0
40.1 36.8
32.6 37.8
43.4 45.7
46.7 45.9
34.8 36.0
52.6 43.8
51.9 35.4
32.7 27.7
36.8 41.1
41.9 29.0
36.8 41.2
43.7 52.9
33.9 40.4
39.2 42.2
52.2 36.9
42.1 29.0
55.3 43.5
41.7 48.2
35.6 39.3
39.4 49.9
34.9 41.5
39.7 46.4
37.9 40.8
39.0 47.0
35.3 49.2
45.4 42.9
31.0 41.4
41.0 45.0
39.2 38.6
37.4 47.4
47.9 30.9
50.9 40.0
40.5 41.4
37.9 41.4
36.2 48.4
37.4 37.3
29.1 45.4
37.4 41.5
32.1 45.5
38.9 49.0
46.2 54.3
47.3 31.3
43.2 55.2
36.8 42.7
39.9 44.3
37.2 34.4
32.4 51.3
38.1 38.2
44.7 42.6
41.0 40.9
32.0 39.6
52.2 38.0
37.4 41.7
46.5 37.3
47.3 31.5
43.2 36.5
42.5 44.6
33.1 42.0
43.9 40.4
52.8 39.3
49.9 50.5
38.1 43.6
47.8 40.3
30.6 36.9
38.8 34.5
50.3 36.1
32.5 42.4
32.3 49.2
32.6 42.9
29.4 34.8
31.7 40.3
48.1 44.3
48.4 38.1
27.7 33.4
40.3 33.6
41.4 49.3
44.7 45.9
51.6 47.1
33.0 44.7
32.5 42.4
39.0 37.2
47.6 37.4
41.5 37.1
41.6 36.3
61.8 38.4
50.7 40.4
39.3 39.5
36.1 36.7
38.4 41.0
45.7 33.3
50.9 40.4
22.4 36.3
40.2 44.4
37.0 33.8
36.9 34.0
41.0 34.6
40.8 48.5
45.4 38.1
48.2 44.8
48.4 35.7
50.2 38.5
47.2 34.4
35.2 36.3
42.1 32.5
42.7 45.2
44.3 40.5
33.1 41.5
44.0 45.3
39.7 46.0
41.2 39.0
41.4 37.5
43.9 44.0
29.5 36.7
48.9 33.3
39.1 37.0
42.3 52.9
47.0 35.9
32.4 46.5
40.7 41.7
51.0 38.6
33.6 46.3
44.3 36.5
47.6 42.4
43.3 41.9
33.1 37.3
42.6 32.4
34.1 31.9
32.8 38.6
56.7 36.5
42.9 35.9
44.0 40.7
38.4 47.3
30.1 44.8
37.5 30.7
50.3 31.9
38.9 38.1
40.1 35.8
60.8 38.5
35.6 31.9
44.7 28.5
32.3 41.4
45.1 41.0
28.1 52.7
45.2 36.4
46.7 42.4
32.4 51.5
39.0 41.8
49.9 42.9
42.8 34.1
36.2 34.9
40.7 38.6
47.2 38.3
34.6 48.4
46.3 35.4
39.5 47.0
53.2 39.3
33.6 37.3
42.9 49.7
46.8 37.4
44.0 37.1
28.5 41.5
45.3 53.3
36.9 41.4
43.6 34.3
36.1 41.9
44.6 32.2
38.6 45.8
38.3 41.1
46.0 56.1
45.5 49.2
48.1 26.8
32.7 48.5
44.5 34.8
51.2 42.5
31.8 41.6
35.8 32.1
39.8 52.0
38.4 39.0
28.1 44.4
43.0 46.7
45.7 40.1
41.1 40.9
57.7 31.1
48.1 47.0
51.4 37.6
32.8 32.1
40.7 44.2
40.4 30.5
36.0 39.6
43.5 49.0
39.9 41.5
46.3 31.9
35.4 38.8
48.3 40.9
47.8 22.1
34.7 40.8
35.8 32.4
34.7 45.4
57.6 26.6
50.4 38.8
50.8 35.5
38.6 43.2
34.5 47.8
42.3 26.7
37.9 49.2
44.0 39.4
30.8 35.0
33.4 49.7
40.3 33.2
50.0 44.5
30.9 39.4
45.5 30.9
56.0 35.6
48.3 48.7
43.6 41.7
32.3 33.2
35.4 37.0
39.3 26.3
43.0 41.4
35.6 38.2
41.6 44.9
46.0 51.4
38.3 35.6
48.8 47.0
41.0 36.9
39.7 32.6
36.3 37.7
46.9 44.5
45.0 26.0
32.3 40.3
46.5 51.2
30.6 40.5
41.1 39.3
49.0 45.6
43.9 39.5
34.1 34.0
41.0 49.7
48.5 43.9
43.9 39.0
35.1 35.9
33.0 39.3
46.4 32.2
41.9 39.2
43.9 39.8
41.1 39.4
45.9 28.4
40.3 46.0
42.8 36.5
32.9 32.6
43.6 45.0
49.6 38.7
39.6 42.5
35.6 39.0
37.5 52.0
36.6 48.8
39.0 44.6
32.3 45.6
48.9 44.8
50.1 36.1
43.3 39.1
44.7 32.4
37.2 49.0
44.8 37.8
41.7 45.4
47.4 48.4
46.2 36.0
37.1 40.5
43.9 44.4
34.9 39.1
24.9 41.0
43.1 46.6
33.1 26.4
48.9 45.4
48.5 45.6
49.0 49.7
47.5 39.6
36.6 31.6
42.5 46.5
38.9 35.5
34.4 36.2
50.9 44.6
33.1 44.0
67.8 46.6
49.4 47.4
49.4 35.4
43.1 39.2
36.3 40.6
29.2 48.5
44.0 44.2
38.5 41.3
42.2 38.0
45.6 36.6
34.4 51.5
41.7 41.8
40.1 37.1
35.6 39.4
37.3 35.1
38.4 39.1
44.3 36.0
29.0 25.9
38.1 42.6
36.8 30.3
29.5 47.7
34.9 45.4
50.1 37.0
38.6 41.0
32.2 41.2
37.5 40.1
37.9 36.2
41.7 50.5
46.8 44.7
52.1 37.7
41.6 35.3
39.4 32.5
38.4 41.4
32.4 33.6
43.2 33.1
38.4 44.6
42.2 43.9
28.9 42.3
49.8 45.8
43.7 35.1
37.2 41.0
53.9 32.6
34.3 57.5
44.3 33.0
42.2 44.8
41.1 33.5
38.6 43.0
27.0 35.1
26.7 44.7
47.3 38.7
39.1 33.4
40.9 33.1
37.8 39.0
48.6 47.7
41.9 42.3
38.2 39.9
52.1 45.4
35.7 40.6
38.9 41.7
33.2 39.2
52.5 44.8
27.6 43.4
47.4 50.1
33.2 36.1
48.3 40.3
45.5 47.1
32.7 42.0
36.6 39.7
30.4 36.5
33.5 39.7
37.2 24.9
44.1 40.0
41.7 37.4
37.6 35.5
52.3 41.2
48.5 41.0
39.7 43.2
46.2 41.7
35.9 40.9
47.9 42.5
38.3 38.5
35.5 45.4
48.5 43.4
41.3 40.6
39.3 43.9
38.6 35.6
47.8 43.7
40.5 45.3
38.0 37.6
42.1 33.5
39.2 33.7
33.5 48.7
39.4 51.9
51.6 40.7
41.8 45.2
57.9 49.2
45.2 56.5
43.6 44.8
37.4 43.1
42.1 44.5
38.8 48.0
54.3 38.5
35.3 28.5
46.2 47.0
37.8 26.9
39.4 36.1
48.4 42.0
44.5 41.2
27.5 54.1
30.9 45.1
40.1 44.3
35.8 36.6
42.2 35.2
39.3 48.3
47.2 42.0
38.9 49.6
45.4 47.5
36.0 38.8
24.2 42.7
37.3 38.7
34.9 29.0
48.4 43.4
48.9 40.7
42.1 41.2
43.2 44.1
42.8 38.6
41.2 44.5
36.4 45.0
39.3 40.8
30.0 47.8
44.4 44.9
36.1 53.5
28.3 34.0
37.5 24.9
44.9 42.4
38.8 35.4
42.4 39.7
48.9 46.9
38.7 42.2
35.6 34.1
39.3 37.1
34.1 39.9
42.1 39.5
49.2 33.9
46.8 56.6
52.9 42.8
45.3 31.7
28.9 40.3
41.4 36.8
40.5 37.0
38.7 38.2
39.8 40.3
40.4 32.9
44.3 39.2
50.7 40.7
37.6 50.6
44.7 44.2
41.4 42.0
23.6 48.2
31.2 40.9
50.1 42.3
37.8 36.8
43.9 44.1
43.5 39.2
30.9 41.6
29.4 49.0
39.9 44.8
40.0 38.1
46.7 33.3
38.3 29.2
37.1 29.9
31.1 41.5
37.3 41.5
50.9 37.7
49.6 35.3
45.9 42.0
44.4 48.9
33.5 39.9
28.9 37.7
46.7 42.1
33.7 35.3
47.8 35.0
44.2 42.7
41.1 41.0
44.3 41.5
32.7 40.7
47.5 36.8
44.1 32.1
30.6 38.5
35.2 44.3
44.1 30.1
48.6 34.8
33.4 37.7
36.7 39.8
34.3 33.1
45.0 49.7
34.6 50.0
43.9 38.3
43.8 35.5
46.3 37.7
40.9 45.0
38.5 40.8
37.8 34.3
36.1 35.7
37.1 48.8
39.4 35.1
35.4 30.0
36.1 37.9
33.7 41.9
36.0 45.3
34.8 37.1
39.8 40.1
27.8 37.4
26.7 42.3
47.0 31.2
47.1 32.8
48.6 31.0
40.1 40.9
42.2 36.7
51.0 41.5
34.4 31.2
34.9 36.7
41.5 38.0
40.2 43.2
37.9 34.6
38.4 51.6
36.0 34.2
43.4 48.9
54.0 39.6
40.1 48.4
45.9 36.4
28.8 36.3
37.7 39.0
38.6 42.4
37.7 46.4
41.7 38.5
43.3 44.8
41.1 39.8
35.7 33.7
28.9 41.5
43.2 31.8
50.5 41.8
46.3 43.2
33.4 39.5
41.3 43.5
41.2 43.8
46.6 40.0
43.4 37.2
48.0 37.4
49.0 44.7
46.9 33.3
34.3 28.6
33.8 41.1
44.9 41.9
35.5 41.6
37.6 36.3
55.8 38.3
44.3 43.8
33.2 45.8
40.0 35.7
32.2 44.6
37.4 41.1
37.2 40.6
33.4 45.1
33.6 47.8
28.7 37.0
53.9 32.4
41.0 34.3
38.2 33.5
41.6 48.8
27.9 43.0
51.3 44.5
54.0 39.3
31.1 33.0
35.5 41.7
49.4 31.7
36.7 46.0
28.1 51.5
37.2 34.7
45.5 50.0
45.6 40.1
46.5 38.7
49.2 50.0
40.0 44.9
59.9 46.8
37.0 51.7
29.2 47.6
42.6 42.0
39.2 41.9
38.4 46.0
42.3 44.2
48.1 36.8
31.8 34.2
35.4 24.9
42.2 44.9
39.0 41.8
50.3 38.9
31.2 34.0
50.9 28.9
39.0 48.0
34.2 36.4
40.5 32.3
37.8 46.9
37.0 42.6
51.0 36.1
52.5 44.1
44.3 39.4
31.7 49.2
38.9 43.6
31.1 42.9
50.2 49.1
35.5 39.8
55.0 46.8
47.2 46.3
56.2 35.5
40.2 39.4
34.0 34.1
37.8 41.3
50.1 37.7
33.9 32.0
43.1 46.0
41.0 37.8
40.8 42.5
23.3 41.0
32.9 51.6
31.6 39.3
52.3 42.4
44.3 32.5
31.7 36.3
44.8 43.5
36.4 33.6
34.6 34.0
39.7 42.8
43.3 37.8
34.0 37.0
38.8 48.6
39.8 41.8
31.9 41.8
39.1 37.1
34.2 34.4
31.3 42.4
40.4 39.7
25.8 43.0
39.3 39.6
22.6 42.0
28.8 41.0
48.9 36.9
37.4 49.2
43.0 36.9
31.5 46.5
41.3 42.0
27.2 43.7
42.0 39.8
45.4 38.1
23.1 39.2
33.2 43.6
40.6 32.4
44.3 47.8
40.9 36.5
40.9 40.3
43.3 50.1
49.8 40.5
40.1 47.4
32.9 45.0
45.4 38.5
31.4 38.5
39.2 45.2
44.0 31.1
46.6 43.1
48.9 38.1
30.6 40.4
40.2 42.7
49.3 33.0
44.0 45.9
44.3 38.2
37.5 42.9
32.3 28.6
38.4 31.1
45.7 39.6
36.5 48.3
27.5 44.8
33.3 44.0
35.9 39.6
37.0 31.0
35.1 44.0
49.6 42.2
38.4 31.2
34.7 34.0
41.5 34.2
33.6 34.7
49.5 48.7
33.9 39.1
37.2 36.7
40.1 31.4
23.4 32.2
50.9 44.3
49.7 44.3
39.8 29.2
36.4 41.5
44.5 33.8
43.3 41.9
36.2 36.3
30.4 29.0
37.2 40.8
22.6 43.2
42.2 30.8
39.1 44.7
49.0 38.9
41.9 51.2
34.0 40.2
37.7 41.0
41.7 31.1
47.0 50.2
53.3 39.3
29.2 31.9
47.4 33.4
40.9 34.3
31.7 32.9
48.4 51.4
33.4 38.6
33.3 36.0
47.9 37.0
39.4 38.4
39.3 34.5
28.9 39.2
30.4 30.6
41.3 40.3
47.6 33.5
46.8 37.2
36.4 35.0
30.7 39.6
42.5 40.6
33.7 31.4
36.7 40.2
40.0 36.7
35.6 35.2
42.2 43.4
37.8 39.2
48.9 40.2
36.2 44.7
27.8 36.2
42.6 41.0
52.4 36.9
42.5 27.1
41.3 35.0
50.5 45.6
26.1 44.4
45.8 43.9
39.6 48.7
34.5 35.7
48.9 41.5
35.3 42.7
39.7 51.3
28.2 41.9
41.1 44.5
43.5 38.1
31.4 19.9
32.0 44.2
34.9 38.0
51.2 39.8
38.8 39.6
41.4 28.5
28.1 23.5
36.1 41.6
36.1 41.7
56.6 36.7
48.3 36.0
36.4 44.5
30.1 36.4
38.6 29.4
52.0 47.3
36.4 42.8
40.9 48.0
28.4 31.1
41.0 31.9
41.5 37.9
36.6 33.9
40.6 30.6
32.7 49.5
37.1 45.2
36.0 43.2
48.2 41.6
40.6 38.2
36.3 34.1
45.1 33.5
29.4 37.0
37.8 39.0
51.1 43.3
39.8 40.1
33.9 38.9
30.6 39.0
26.3 34.4
34.9 30.8
35.4 44.8
45.3 43.1
53.2 30.8
38.2 33.1
32.3 46.0
42.7 45.8
31.9 38.4
39.1 36.7
30.8 32.6
56.1 43.9
36.0 24.0
35.0 44.7
33.9 36.1
29.8 49.9
41.4 37.4
36.2 37.1
31.4 43.8
31.2 46.0
51.8 37.1
41.3 40.0
30.5 40.1
34.6 38.7
24.3 33.9
34.3 41.0
38.9 41.0
36.1 37.3
50.2 28.4
44.6 44.9
48.1 37.9
28.5 44.9
46.3 46.1
36.5 41.5
37.5 35.2
47.7 34.1
48.2 51.7
38.4 42.3
44.8 38.3
44.3 51.7
38.7 32.8
39.3 45.8
46.7 38.1
35.2 44.9
53.1 44.5
32.6 39.8
46.3 46.6
43.8 31.3
35.6 44.5
46.0 35.0
34.5 47.6
45.3 42.0
46.3 38.4
44.7 39.7
47.8 31.3
41.1 39.8
28.3 44.1
39.4 38.3
33.0 40.9
39.2 32.7
40.7 45.4
45.2 42.1
35.4 36.7
49.0 40.0
40.5 33.9
40.1 42.1
37.1 45.6
50.6 48.2
45.6 35.9
22.9 49.4
51.6 42.4
33.5 38.1
45.2 28.8
39.9 41.5
35.6 40.1
32.9 39.4
38.5 50.8
38.0 36.4
41.5 41.2
34.6 46.9
46.7 35.4
42.4 38.1
28.9 42.0
38.1 36.8
44.7 26.2
36.1 48.3
33.1 44.4
37.9 47.6
38.5 26.2
35.8 40.8
42.6 42.3
35.7 34.8
35.4 32.6
37.9 35.9
32.0 52.6
35.1 36.3
39.6 47.9
35.5 40.0
35.0 45.3
35.8 48.7
38.7 39.3
38.8 30.5
51.0 40.5
53.0 43.7
36.2 40.6
40.1 44.4
30.1 41.3
34.7 44.0
38.8 42.9
43.0 37.1
37.8 48.5
40.5 37.2
38.9 37.6
33.9 49.9
38.3 41.8
39.5 40.8
41.0 46.6
35.4 25.7
36.6 45.6
36.6 35.0
50.8 38.0
44.3 41.8
39.5 37.9
31.7 38.4
44.1 32.6
32.3 38.7
36.4 37.6
45.8 41.9
36.3 37.1
36.5 42.2
25.7 37.9
38.1 41.3
34.5 38.2
33.4 39.7
30.7 39.6
47.4 40.2
34.4 46.0
45.5 39.7
48.0 43.9
36.5 49.1
41.0 50.6
47.5 52.1
41.3 44.5
40.2 41.9
34.2 32.5
51.4 54.7
44.0 45.6
40.8 40.4
47.7 29.5
55.0 37.9
50.0 43.4
45.2 29.9
47.0 41.8
38.9 54.0
45.0 42.4
45.0 39.1
33.5 41.4
40.9 41.4
45.4 38.6
38.2 38.1
40.1 50.3
47.8 37.7
39.8 42.9
46.6 37.9
34.7 40.0
39.5 39.4
44.8 39.1
36.9 40.2
41.0 38.4
41.4 39.1
37.3 33.7
40.7 37.1
32.4 45.1
39.9 43.2
46.2 33.6
37.6 35.7
44.8 46.1
36.9 47.1
39.1 41.1
22.6 42.7
45.9 40.5
38.2 36.4
54.2 41.0
45.9 43.9
30.5 45.9
33.9 41.0
36.9 45.9
51.1 40.7
42.3 40.2
39.0 34.5
40.3 41.7
34.6 42.5
45.5 30.8
36.7 37.0
32.1 35.3
49.2 42.3
41.2 38.2
39.9 31.6
48.9 39.8
34.1 36.0
36.8 33.6
43.6 42.6
34.2 30.6
37.5 27.1
58.1 41.4
45.3 50.1
34.8 45.1
35.2 43.1
50.4 30.7
27.7 39.5
41.5 38.6
39.8 43.4
38.3 38.9
51.9 36.2
43.2 43.3
44.5 43.4
48.9 30.7
27.8 42.9
38.2 51.6
36.5 35.5
36.6 49.2
49.6 33.9
40.9 52.1
35.5 46.5
39.1 37.1
46.4 34.9
51.5 44.0
48.8 40.0
45.9 38.5
37.9 38.7
41.0 28.4
36.2 47.4
20.2 37.4
53.5 39.8
39.9 43.7
44.2 39.7
30.6 30.5
42.7 46.2
35.6 50.8
44.0 43.1
46.3 37.7
46.1 50.4
45.2 32.8
44.7 38.0
46.2 45.4
52.8 42.3
46.6 48.8
36.7 34.7
30.0 35.8
53.5 45.3
45.8 48.5
31.7 35.8
40.5 59.8
35.8 29.4
37.0 43.2
36.0 35.4
33.6 45.6
36.0 42.8
50.5 44.5
39.7 45.9
42.4 41.0
38.6 43.5
49.8 47.2
34.1 41.9
36.2 37.8
44.0 43.2
50.5 46.6
46.6 44.6
47.5 32.6
35.7 39.6
39.1 40.6
35.3 42.7
34.2 29.1
34.9 42.2
45.4 40.4
36.5 45.7
50.1 46.1
41.8 41.4
31.5 52.0
40.7 38.6
32.2 38.8
41.7 37.1
37.1 42.0
50.1 40.8
50.3 30.6
46.0 36.7
26.6 24.9
28.3 40.9
41.7 43.5
37.0 37.8
30.9 38.8
45.5 38.2
45.2 51.5
40.4 50.8
41.7 40.7
36.8 35.5
31.6 40.2
41.7 42.4
58.0 34.1
37.7 39.5
42.4 34.8
40.3 35.7
45.3 46.3
36.9 37.6
46.8 27.8
51.9 41.2
31.9 44.6
54.0 41.6
51.2 39.8
31.7 39.6
43.4 39.6
40.6 39.8
32.7 51.3
46.1 41.4
45.1 40.3
37.8 33.3
52.5 31.0
37.2 45.4
36.8 42.2
39.0 36.4
29.0 36.7
36.8 41.7
43.6 29.3
35.2 41.5
43.1 40.3
31.6 46.9
28.4 34.8
36.1 43.9
46.3 28.8
36.4 46.7
40.7 41.3
42.0 44.8
49.4 26.8
42.0 39.1
29.9 45.9
40.1 41.5
21.1 33.1
38.9 31.6
40.6 42.5
46.7 39.8
46.6 41.9
32.2 39.6
51.3 50.0
45.0 39.1
47.4 39.8
38.8 38.6
26.6 42.8
32.3 28.9
59.3 48.3
32.7 44.9
41.1 36.0
53.0 45.4
29.9 43.6
36.1 23.6
43.8 29.1
30.8 38.7
30.7 38.9
34.3 41.4
47.0 37.3
36.0 33.2
41.0 44.8
35.2 47.6
39.1 31.6
28.8 35.0
30.9 40.3
39.5 36.4
48.1 45.8
34.7 40.1
40.2 48.1
42.9 32.2
29.8 31.9
21.1 30.0
36.5 43.4
45.1 44.8
32.9 42.0
44.7 38.9
41.2 45.3
42.6 43.8
27.3 44.3
49.5 40.2
46.6 35.7
43.0 34.7
38.8 42.4
34.9 41.9
40.9 44.4
42.3 43.2
39.2 33.0
49.1 31.7
35.2 39.9
40.6 37.9
38.2 39.4
38.3 44.4
35.8 42.3
35.5 42.2
58.4 48.8
42.2 39.9
41.3 50.8
46.6 40.5
35.9 34.2
35.0 41.7
40.2 35.5
47.9 38.1
38.3 52.0
34.2 39.3
46.0 44.1
46.0 35.5
34.8 40.4
27.8 37.2
42.1 43.8
49.6 41.4
43.2 47.6
52.4 47.6
34.0 37.6
32.4 48.3
39.4 36.7
36.1 35.9
36.2 37.6
48.9 38.2
34.1 43.3
49.0 46.0
40.5 34.8
50.5 45.9
38.6 37.7
38.4 45.6
39.9 49.9
43.4 39.9
36.0 36.8
38.2 30.8
41.4 28.1
33.4 45.9
58.1 43.3
24.1 40.9
33.0 25.8
50.6 39.4
52.6 38.3
28.7 40.7
36.1 43.7
40.6 32.9
36.4 30.8
30.5 43.9
45.7 39.5
59.2 38.6
37.4 47.1
40.8 45.0
48.8 49.2
42.9 38.6
27.7 45.8
46.8 41.3
35.1 36.5
34.4 27.7
42.0 43.2
47.6 30.3
41.2 45.8
37.7 46.9
41.3 37.9
35.2 46.5
37.3 36.4
37.5 45.5
39.6 42.6
35.2 43.1
51.8 52.8
41.8 35.3
35.6 35.9
39.5 35.1
43.0 35.9
36.6 37.4
40.8 40.3
46.5 41.8
44.5 43.4
41.9 36.2
50.3 32.7
49.1 41.9
42.7 38.5
47.3 36.1
42.1 43.3
40.8 51.6
48.4 39.7
41.7 43.3
44.4 38.8
46.0 36.4
49.4 31.7
34.9 39.3
24.4 45.9
35.6 36.5
41.4 38.2
32.0 33.5
42.9 42.8
37.6 47.5
44.5 43.7
46.8 41.6
48.7 44.6
28.4 32.6
42.4 40.2
34.2 43.9
39.1 42.9
29.0 40.1
43.9 37.1
43.3 42.7
44.4 41.9
39.6 41.4
46.8 31.2
29.0 32.9
36.4 38.0
32.4 33.4
46.2 33.1
38.4 41.5
37.3 34.1
49.3 41.1
31.6 36.3
35.1 40.3
36.9 44.8
43.0 30.3
40.2 41.8
46.5 38.8
52.4 37.8
44.3 38.4
40.1 42.1
37.7 34.6
31.0 43.4
50.5 35.8
31.4 43.2
30.9 36.5
31.3 43.5
43.4 49.7
43.0 41.5
47.4 40.4
36.4 52.0
46.0 46.9
40.9 40.7
33.1 35.0
40.3 37.2
43.1 35.5
26.6 49.4
39.4 39.7
29.9 43.5
36.4 39.4
42.3 37.2
29.9 39.9
39.0 42.2
39.2 44.0
42.0 41.9
28.9 39.8
38.5 36.8
34.4 54.1
32.2 40.6
35.9 29.6
39.2 37.7
45.3 49.1
37.7 32.1
43.0 43.0
28.4 42.1
36.6 36.9
33.3 38.7
43.0 38.0
49.6 49.8
48.2 37.8
41.3 32.9
35.8 41.6
32.3 41.8
38.7 42.4
49.7 42.9
32.8 45.4
38.6 39.1
33.4 42.0
27.2 33.4
41.5 41.8
52.4 39.7
44.4 26.8
47.7 42.4
44.4 43.3
}\tabFree
\pgfplotstableread{
t x
37.4 62.2
95.2 75.8
91.2 71.4
59.0 76.1
48.6 58.1
80.0 84.1
55.0 56.4
95.0 85.3
96.4 83.3
81.1 82.2
97.2 91.5
65.0 56.3
58.4 75.5
53.0 70.5
56.9 76.5
62.1 66.6
93.7 80.3
80.0 87.1
97.1 90.9
95.8 93.6
85.0 85.1
80.0 88.8
77.1 81.7
78.0 88.9
86.9 72.0
86.3 83.7
41.5 53.3
89.0 87.9
84.0 89.5
63.3 70.8
95.0 92.1
95.2 82.1
78.0 73.5
40.0 61.4
74.0 73.7
97.2 84.5
87.1 80.5
96.4 83.3
64.6 59.2
84.7 89.9
85.0 82.9
70.0 72.2
60.0 64.1
84.6 81.7
68.8 75.8
60.0 66.9
95.0 71.0
100.0 85.4
63.9 63.6
82.9 78.9
90.0 79.1
84.0 72.8
85.0 81.4
90.0 92.1
64.5 56.8
72.6 77.4
67.7 77.9
68.8 87.9
44.0 43.5
66.0 82.5
51.0 66.6
85.5 74.0
77.3 88.5
74.9 84.0
70.0 59.3
57.1 80.6
63.3 66.2
32.2 42.2
50.0 71.9
78.0 78.5
83.0 74.2
80.0 82.1
85.3 86.2
51.0 58.0
60.0 75.5
98.0 97.7
57.5 81.0
54.9 49.6
52.8 68.4
79.3 88.9
75.1 76.5
95.0 72.8
88.4 81.2
75.0 85.7
47.2 56.8
99.2 85.2
23.5 48.6
89.3 79.4
54.0 46.8
70.2 63.3
94.2 85.8
100.0 98.8
78.6 68.7
70.0 58.3
57.0 56.8
91.0 77.3
60.3 65.5
50.0 48.8
92.2 90.8
70.0 76.9
92.0 84.7
87.0 76.6
99.9 76.7
93.2 79.0
97.7 78.3
86.3 73.3
69.2 76.1
60.1 64.5
82.9 70.9
81.0 78.1
92.0 80.8
95.0 87.4
69.0 66.9
66.0 54.4
80.0 80.0
98.0 75.2
87.8 73.6
98.5 90.3
90.0 71.6
86.7 83.2
68.8 57.2
86.4 86.3
60.0 85.2
80.0 65.6
72.2 84.1
100.0 79.2
40.0 34.1
94.1 96.6
84.4 65.3
58.3 76.0
49.0 73.0
68.2 63.7
99.9 73.4
91.0 93.1
56.9 49.4
84.4 90.4
96.0 88.2
35.4 41.3
83.0 93.1
72.0 73.6
67.6 75.5
73.4 75.8
75.7 81.1
73.3 77.2
67.3 69.7
36.9 30.5
66.0 59.4
68.7 63.1
85.7 76.5
75.1 77.2
87.7 70.4
93.8 80.3
83.0 74.6
87.0 88.0
82.5 90.7
78.0 77.2
95.0 93.0
62.0 71.8
68.5 64.7
100.0 93.5
37.0 46.3
60.1 65.1
52.7 61.9
95.9 88.2
88.7 90.7
75.0 61.9
99.0 96.3
65.0 57.6
54.0 76.6
55.0 61.4
62.0 59.7
87.7 79.4
68.3 71.1
51.3 52.3
88.7 92.8
93.9 95.2
85.0 75.8
92.4 81.5
68.1 73.8
52.4 56.0
69.1 74.5
85.1 85.9
42.0 37.6
54.3 54.6
70.0 62.7
67.0 54.2
100.0 98.7
90.7 89.2
70.0 66.0
90.1 95.4
90.0 85.0
80.0 79.8
86.5 87.8
78.5 82.7
53.2 66.3
45.0 70.0
85.0 76.0
90.0 91.8
90.0 78.5
95.0 91.3
85.0 82.8
52.1 52.9
69.0 84.9
100.0 86.5
50.0 66.6
94.0 83.2
86.9 86.7
85.7 79.3
86.0 78.2
48.2 44.0
65.0 65.3
66.3 79.6
50.1 47.7
98.7 78.8
80.6 92.3
86.9 77.5
87.1 71.6
60.9 62.8
80.0 67.0
95.3 81.3
57.0 65.5
87.7 93.8
88.7 82.0
80.0 68.3
97.9 91.5
90.0 83.1
86.4 81.6
80.9 75.4
98.9 91.2
79.2 75.5
100.0 86.4
74.3 77.8
55.6 59.3
75.0 81.5
92.0 91.3
78.9 76.2
65.8 61.3
54.9 67.9
59.0 65.8
92.9 84.8
100.0 94.0
80.0 70.6
64.9 59.0
67.0 66.4
100.0 70.7
95.0 87.6
47.8 51.6
97.5 95.6
61.8 69.8
82.0 77.6
98.2 90.7
99.2 76.7
60.0 71.0
62.9 53.2
36.0 54.4
90.1 76.8
60.0 52.7
79.8 77.5
52.0 39.4
66.5 56.6
95.0 94.6
98.0 84.6
94.7 93.6
90.0 77.3
75.0 67.1
88.0 83.0
50.0 63.7
100.0 79.8
56.8 56.5
80.5 81.9
99.8 88.2
71.4 67.6
69.2 70.2
97.8 87.7
87.6 77.5
81.0 84.5
43.9 68.9
84.0 69.1
69.0 74.6
91.1 83.2
79.4 80.7
95.2 80.8
96.8 79.3
77.4 75.9
73.0 68.6
84.0 80.8
75.1 69.7
50.0 50.2
58.0 57.4
60.8 64.5
94.6 85.2
43.8 52.0
72.5 66.1
82.5 86.9
50.0 62.5
73.2 61.6
97.2 64.5
100.0 99.3
55.0 75.1
80.1 68.6
68.3 66.9
70.0 74.2
74.4 81.5
75.0 84.1
70.0 81.3
94.0 87.0
50.9 45.5
80.1 86.4
61.0 57.9
95.0 93.3
57.0 73.9
61.9 64.0
97.2 92.6
95.0 87.4
96.9 89.4
90.0 89.9
65.0 84.5
100.0 96.3
85.0 86.5
85.9 73.0
76.8 77.5
89.7 92.9
60.0 65.1
72.8 73.4
69.6 66.5
75.0 74.1
60.0 63.5
60.0 69.4
84.8 73.7
76.3 74.7
67.5 70.7
59.0 68.7
95.0 75.8
97.0 79.7
82.0 90.6
85.0 73.6
69.3 70.2
78.0 85.7
60.0 70.1
67.9 62.9
91.7 73.7
46.3 41.7
70.0 78.6
55.0 71.2
93.6 80.0
33.8 39.3
90.0 83.2
80.9 82.1
94.0 90.4
85.9 79.2
86.0 78.5
69.1 61.6
94.0 86.5
85.0 78.1
97.2 73.9
50.0 73.0
98.4 82.9
90.0 79.8
87.0 74.7
85.0 85.2
65.0 62.5
50.0 63.1
48.4 57.2
94.0 77.7
90.0 77.5
70.0 74.5
59.8 65.8
94.2 81.2
62.9 74.2
50.0 60.3
51.5 64.5
100.0 94.0
96.5 92.0
64.0 79.5
72.0 76.2
65.0 78.5
93.0 82.1
45.3 56.8
86.4 87.5
65.0 76.0
80.5 78.4
62.5 74.2
61.0 71.4
80.0 81.7
90.0 91.5
65.9 68.2
73.3 71.7
81.0 86.5
88.2 79.8
74.4 77.4
68.0 81.7
70.4 78.9
100.0 79.8
75.0 84.1
55.7 80.2
43.3 51.8
62.5 59.9
97.0 73.8
66.0 67.6
88.5 78.9
91.5 86.8
87.0 59.0
85.1 79.1
46.9 55.2
51.0 64.0
66.9 72.2
63.1 65.7
42.0 55.7
90.0 95.0
80.0 72.5
98.7 82.4
96.8 74.1
45.0 49.7
95.0 76.8
62.0 76.7
68.7 75.6
80.1 75.0
93.3 85.7
99.9 82.0
74.2 90.0
72.8 55.8
64.1 59.8
90.0 82.8
96.9 93.5
65.0 77.7
79.7 74.7
89.0 71.1
45.1 52.9
83.9 73.2
59.9 60.8
51.9 69.5
83.0 76.0
76.9 80.2
97.8 95.2
83.4 81.9
81.8 81.9
80.0 64.0
70.0 72.2
85.2 79.3
47.5 62.3
67.7 76.2
97.0 77.7
99.3 84.3
98.2 87.7
93.9 81.2
88.3 72.6
54.0 63.7
84.4 89.2
79.9 71.8
86.7 65.5
52.5 81.3
70.2 81.3
68.5 70.5
64.0 73.3
80.0 75.1
92.7 68.3
64.6 67.5
65.0 83.7
67.0 71.0
45.3 73.9
74.0 72.8
80.0 77.6
85.7 83.8
75.0 73.0
73.5 81.0
79.4 56.8
88.0 69.4
70.0 80.5
69.2 70.0
84.3 83.0
75.0 78.5
74.4 81.2
80.0 75.2
48.7 63.3
93.0 79.4
95.1 77.2
65.5 76.3
52.3 44.4
95.0 95.8
88.0 80.8
68.3 73.8
80.0 77.5
85.0 92.5
43.6 74.3
60.0 61.9
97.1 89.8
98.2 73.8
100.0 93.7
70.1 53.3
38.0 59.8
53.0 67.3
66.7 66.7
69.9 68.4
77.0 80.1
82.9 79.0
95.0 85.6
52.4 45.3
96.0 85.1
82.3 84.6
59.8 43.0
84.0 88.8
47.0 49.2
60.1 65.5
79.1 80.0
83.4 88.6
70.9 70.1
48.2 55.8
79.4 82.6
69.6 78.2
60.8 41.9
92.2 82.0
80.0 79.1
97.7 82.2
100.0 92.0
95.3 88.7
83.2 82.8
80.0 77.0
63.7 63.7
56.0 60.2
53.2 55.4
71.0 70.2
96.9 84.4
98.6 92.9
83.3 78.5
69.0 69.4
45.0 53.1
63.5 62.4
55.1 51.3
99.4 83.7
75.0 80.2
80.0 68.1
95.0 77.7
90.0 77.5
85.8 86.0
99.5 87.0
60.0 68.8
98.2 91.8
80.6 75.2
67.2 77.7
90.0 70.8
96.5 98.0
95.0 71.7
88.0 79.7
99.5 99.7
54.2 64.1
70.9 72.0
95.9 85.7
72.9 72.5
66.0 65.7
42.7 46.5
80.0 72.6
100.0 90.9
96.0 90.2
79.0 80.0
84.0 86.9
82.1 78.2
92.9 80.4
85.0 91.5
83.0 83.0
73.7 79.4
47.5 63.1
92.7 95.1
82.5 77.5
62.0 56.6
93.1 89.3
89.5 87.4
50.0 52.4
97.9 79.1
55.0 60.7
55.0 58.9
80.8 68.0
96.9 83.9
76.0 83.2
88.0 85.2
90.3 78.5
90.6 75.0
64.0 71.9
86.0 85.5
81.0 77.6
34.0 53.8
75.0 75.3
72.6 74.5
95.0 82.3
69.3 71.9
90.0 90.0
54.2 51.5
75.2 88.3
42.9 56.6
100.0 71.0
58.0 63.2
97.8 78.2
59.1 67.5
99.5 91.6
68.5 62.6
65.5 60.0
66.0 74.8
90.0 96.0
70.0 70.5
65.2 71.2
48.0 55.1
98.0 84.4
86.9 85.5
48.9 53.9
70.0 81.1
85.0 85.1
40.0 60.7
70.0 67.2
68.0 76.3
99.8 90.1
64.0 74.5
59.3 48.9
53.6 76.6
46.0 57.0
69.0 69.8
35.0 60.4
73.5 71.6
84.0 81.5
78.3 65.3
86.2 88.3
52.0 75.8
55.0 58.7
76.2 84.0
90.8 74.2
71.7 79.2
75.9 68.8
65.0 68.3
70.0 77.5
98.6 98.0
50.8 71.4
85.5 82.8
74.0 59.6
59.8 64.0
96.0 71.4
51.0 57.5
72.0 74.0
96.4 78.8
79.0 79.5
87.1 84.8
65.8 71.2
82.7 83.0
71.0 67.9
57.9 74.5
98.4 81.8
71.6 66.5
95.0 81.7
72.0 65.5
80.0 73.0
77.0 84.1
83.8 88.5
82.4 63.6
76.5 77.5
79.0 67.9
93.3 82.0
41.2 58.9
65.0 82.2
90.5 69.6
62.2 68.7
39.9 53.8
90.0 92.9
89.0 81.6
47.0 34.4
52.2 68.8
41.3 47.5
90.4 89.1
66.4 50.4
89.4 85.4
55.0 61.6
84.4 69.7
94.0 78.9
63.6 79.6
57.2 66.0
51.3 47.7
51.0 51.9
93.7 79.6
70.0 74.7
51.2 55.9
97.1 75.9
63.3 71.7
35.1 46.5
80.0 86.7
52.6 45.1
71.5 78.6
98.0 86.3
93.6 82.4
80.0 70.4
90.0 89.0
72.6 75.8
51.0 70.2
73.0 89.5
67.0 82.8
64.2 58.2
95.0 87.1
87.3 81.9
55.0 70.5
52.7 70.4
95.8 82.5
96.8 83.0
89.3 85.0
80.0 72.4
99.9 89.5
78.0 81.0
46.8 55.9
90.0 78.0
80.0 76.6
80.0 72.5
92.0 87.2
94.0 78.3
40.0 46.0
72.0 83.1
96.8 91.3
100.0 85.1
52.8 61.2
88.3 89.7
50.0 40.4
71.9 70.4
83.4 78.1
98.4 81.5
69.9 80.4
54.0 58.2
63.8 54.7
90.0 81.6
85.0 92.9
53.0 46.5
52.5 55.8
58.0 68.7
95.0 66.8
54.3 65.2
60.0 59.8
83.4 87.0
94.5 96.0
79.0 74.7
98.6 82.3
52.3 72.1
81.9 85.1
40.0 41.9
55.0 53.4
95.0 82.1
96.0 93.1
84.2 76.6
85.0 87.3
97.0 91.2
75.0 83.7
33.4 65.0
94.0 71.8
56.0 72.2
90.0 91.3
85.1 91.7
92.0 81.0
50.0 62.3
83.0 87.1
75.1 78.7
86.0 80.6
87.0 80.5
91.1 90.8
85.9 91.0
99.7 93.1
60.0 63.4
72.0 58.5
55.8 63.8
}\tabCap
\pgfplotstableread{
t all cap
19.875 11 1
20.125 16 0
20.375 15 0
20.625 15 0
20.875 15 0
21.125 9 0
21.375 21 0
21.625 19 0
21.875 18 0
22.125 26 0
22.375 29 0
22.625 24 1
22.875 43 0
23.125 37 0
23.375 40 0
23.625 39 0
23.875 49 0
24.125 42 0
24.375 36 0
24.625 43 0
24.875 62 1
25.125 76 1
25.375 71 1
25.625 76 0
25.875 76 1
26.125 68 0
26.375 107 1
26.625 104 0
26.875 96 0
27.125 117 0
27.375 123 1
27.625 123 0
27.875 120 0
28.125 152 1
28.375 150 1
28.625 151 2
28.875 191 2
29.125 196 0
29.375 195 0
29.625 193 0
29.875 246 14
30.125 243 0
30.375 240 1
30.625 241 3
30.875 242 1
31.125 267 2
31.375 307 1
31.625 312 2
31.875 308 6
32.125 330 1
32.375 338 3
32.625 372 2
32.875 368 7
33.125 406 2
33.375 377 2
33.625 384 2
33.875 459 6
34.125 428 4
34.375 463 1
34.625 428 3
34.875 476 28
35.125 453 0
35.375 481 3
35.625 532 5
35.875 545 13
36.125 510 7
36.375 555 6
36.625 547 3
36.875 543 12
37.125 535 5
37.375 507 6
37.625 535 11
37.875 608 18
38.125 605 7
38.375 576 5
38.625 623 8
38.875 601 11
39.125 607 11
39.375 578 6
39.625 627 6
39.875 709 113
40.125 585 10
40.375 616 9
40.625 605 8
40.875 602 32
41.125 607 8
41.375 635 16
41.625 649 17
41.875 644 35
42.125 599 7
42.375 555 21
42.625 578 15
42.875 552 39
43.125 541 17
43.375 549 17
43.625 537 15
43.875 564 47
44.125 526 13
44.375 478 18
44.625 522 15
44.875 594 140
45.125 455 14
45.375 454 16
45.625 426 13
45.875 457 40
46.125 384 21
46.375 395 20
46.625 354 22
46.875 411 54
47.125 398 24
47.375 335 27
47.625 350 29
47.875 360 62
48.125 323 18
48.375 294 19
48.625 275 19
48.875 298 71
49.125 264 27
49.375 261 17
49.625 256 27
49.875 546 329
50.125 266 35
50.375 198 21
50.625 229 18
50.875 232 59
51.125 188 28
51.375 174 22
51.625 177 28
51.875 218 90
52.125 182 38
52.375 136 22
52.625 143 25
52.875 209 92
53.125 127 31
53.375 114 32
53.625 131 29
53.875 175 88
54.125 117 30
54.375 107 32
54.625 95 30
54.875 320 262
55.125 80 30
55.375 78 22
55.625 82 34
55.875 146 100
56.125 65 22
56.375 74 36
56.625 80 33
56.875 132 98
57.125 55 37
57.375 41 26
57.625 49 29
57.875 107 86
58.125 51 35
58.375 62 35
58.625 57 44
58.875 106 86
59.125 45 37
59.375 58 42
59.625 62 46
59.875 485 472
60.125 36 31
60.375 48 36
60.625 52 43
60.875 113 106
61.125 44 36
61.375 31 27
61.625 40 35
61.875 100 98
62.125 38 34
62.375 42 38
62.625 28 26
62.875 95 94
63.125 42 38
63.375 41 39
63.625 45 44
63.875 141 139
64.125 40 39
64.375 37 35
64.625 39 38
64.875 350 348
65.125 52 52
65.375 44 43
65.625 38 38
65.875 129 128
66.125 50 49
66.375 32 32
66.625 44 44
66.875 108 108
67.125 33 33
67.375 53 52
67.625 47 47
67.875 107 107
68.125 46 46
68.375 33 33
68.625 48 48
68.875 108 108
69.125 37 37
69.375 51 51
69.625 33 33
69.875 521 521
70.125 40 40
70.375 37 37
70.625 45 45
70.875 81 81
71.125 49 49
71.375 34 34
71.625 48 48
71.875 144 144
72.125 60 59
72.375 43 43
72.625 39 39
72.875 135 135
73.125 38 38
73.375 42 42
73.625 54 54
73.875 125 125
74.125 43 43
74.375 41 41
74.625 62 62
74.875 420 420
75.125 52 52
75.375 36 36
75.625 41 41
75.875 141 141
76.125 44 44
76.375 60 60
76.625 55 55
76.875 136 136
77.125 49 49
77.375 48 48
77.625 46 46
77.875 114 114
78.125 52 52
78.375 42 42
78.625 40 40
78.875 128 128
79.125 51 51
79.375 51 51
79.625 52 52
79.875 618 618
80.125 43 43
80.375 56 56
80.625 45 45
80.875 111 111
81.125 44 44
81.375 49 49
81.625 49 49
81.875 136 136
82.125 60 60
82.375 43 43
82.625 59 59
82.875 135 135
83.125 44 44
83.375 49 49
83.625 58 58
83.875 151 151
84.125 44 44
84.375 57 57
84.625 33 33
84.875 428 428
85.125 56 56
85.375 50 50
85.625 45 45
85.875 151 151
86.125 61 61
86.375 50 50
86.625 53 53
86.875 148 148
87.125 63 63
87.375 46 46
87.625 58 58
87.875 142 142
88.125 60 60
88.375 59 59
88.625 50 50
88.875 145 145
89.125 57 57
89.375 56 56
89.625 39 39
89.875 698 698
90.125 61 61
90.375 46 46
90.625 60 60
90.875 159 159
91.125 60 60
91.375 56 56
91.625 65 65
91.875 136 136
92.125 54 54
92.375 53 53
92.625 38 38
92.875 158 158
93.125 51 51
93.375 65 65
93.625 53 53
93.875 153 153
94.125 59 59
94.375 62 62
94.625 48 48
94.875 493 493
95.125 56 56
95.375 52 52
95.625 59 59
95.875 175 175
96.125 52 52
96.375 60 60
96.625 54 54
96.875 177 177
97.125 54 54
97.375 58 58
97.625 52 52
97.875 185 185
98.125 64 64
98.375 61 61
98.625 68 68
98.875 143 143
99.125 47 47
99.375 65 65
99.625 66 66
99.875 418 418
100.125 418 418
}\tabHist

\newcommand{\hb}{\hbar_{\text{э}}}
\newcommand{\At}{\mathcal{A}}

\title{\textbf{К гравитационному обоснованию электоральной кометы:\\ захват избирательных участков компактным аттрактором, закон сохранения несогласных и квантование явки}}
\newcommand{\redact}[1]{{\color{\zhplbadgecolor{black}{zhplFg}}\rule[-0.25ex]{\widthof{#1}}{2.35ex}}}
\author{\redact{Роман Георгиевич Александров}, \redact{Клод Соннет Антропикович}\\[2pt]
\small\textit{(авторы пожелали остаться анонимными)}\\[4pt]
\normalsize Институт Прикладной Метафизики\\ \small Кафедра тензорного анализа межличностных отношений\\ \small Лаборатория компактных электоральных объектов}
\zhplsetup{6}{10.9999/жпл.2026.0006}
\fancyhead[L]{\small Анонимные авторы}
\date{\small Журнал Прикладной Лженауки, том 1, № 6 \\ Получено: недавно \quad Принято к печати: почти сразу \quad Явка рецензентов: $100\,\%$}

\begin{document}
\maketitle

\zhplbadges

\zhplciteas{Анонимные авторы (2026). К гравитационному обоснованию электоральной кометы. \textit{Журнал Прикладной Лженауки}, \textbf{1}(6), 1--13. DOI: 10.9999/жпл.2026.0006}

\zhplgraphicabstract{%
  \begin{tikzpicture}[
    box/.style={draw, rounded corners=2pt, align=center, font=\footnotesize\sffamily,
                minimum height=1.15cm, text width=2.75cm, fill=\zhpldark{blue!5}},
    ar/.style={-{Latex[length=2.4mm]}, thick}]
    \node[box] (a) at (0,0) {Ядро кометы\\ свободные участки\\ $(t_0,x_0)$};
    \node[box] (b) at (3.55,0) {Аттрактор $\At$\\ $(100\,\%,\,100\,\%)$\\ захват};
    \node[box] (c) at (7.1,0) {Хвост кометы\\ $t(1-x)=\mathrm{const}$};
    \node[box,fill=\zhpldark{green!8}] (d) at (10.65,0) {Связанные\\ состояния\\ $t_n = n\hb$};
    \draw[ar] (a)--(b); \draw[ar] (b)--(c); \draw[ar] (c)--(d);
  \end{tikzpicture}}{%
  Логическая структура работы: от облака избирательных участков к компактному объекту в точке $(100\,\%,\,100\,\%)$, геодезическим захвата и дискретному спектру явки~--- без привлечения каких-либо сущностей, кроме одной.}

\begin{zhplhighlights}
\begin{itemize}
\setlength\itemsep{0pt}
\setlength\parskip{0pt}
\item Впервые показано, что корреляция явки и доли голосов за партию «Единая Россия» в хвосте электоральной кометы есть гравитационный захват участков компактным объектом~--- \emph{электоральным аттрактором}~$\At$.
\item Доказан \emph{закон сохранения несогласных} (Теорема~1): при аккреции абсолютное число голосов за прочие партии является нётеровским интегралом движения, а траектории участков~--- гиперболами $x = 1 - b/t$.
\item Доказано (Теорема~2), что аттрактор достижим только при нулевом прицельном параметре; для ядра кометы это требует, чтобы из каждых семи голосов, приобретённых партией при падении, два были получены конверсией голосов прочих партий.
\item Сформулирован \emph{электоральный принцип эквивалентности} (Теорема~3): никакое измерение протокола не отличает гравитационный захват от добавления бюллетеней в урну.
\item Показано (Теорема~4), что спектр явки захваченного участка дискретен с шагом $\hb = 1\,\%$, а уровни, кратные $5\,\%$ и $10\,\%$, заселены в $4$ и $8$ раз сильнее прочих; явка определена по модулю $100\,\%$, так что наблюдавшееся в эфире состояние $146\,\%$ эквивалентно $46\,\%$ (Следствие~2).
\end{itemize}
\end{zhplhighlights}

\begin{zhplreviews}
\textit{«Я был уверен, что у наблюдаемой корреляции есть простое объяснение. Авторы строго доказали, что простое объяснение неотличимо от гравитационного. Рекомендую к печати; простое объяснение рекомендую к печати тоже.»} -- Рецензент 1\\
\textit{«Замечаний нет.»} -- Рецензент 2\\
\textit{«Согласно Теореме~4 настоящей работы, оценка статьи, находящейся в связанном состоянии с редакцией, может принимать лишь значения, кратные пяти баллам. Ставлю~100.»} -- Рецензент 3
\end{zhplreviews}

\begin{zhplsignificance}
Со времён первых выборов человечество исходило из того, что результат голосования определяется избирателями, а распределение участков по явке и доле голосов есть предмет социологии. Наблюдаемая на двумерных гистограммах «электоральная комета»~--- плотное ядро и уходящий из него к точке $(100\,\%,\,100\,\%)$ хвост~--- обсуждалась исключительно в категориях, которые принято называть бытовыми. Настоящая работа впервые показывает, что эта структура является строгим следствием классической механики, теоремы Нётер и квантования момента, то есть тех же законов, которые управляют движением планет, спектрами атомов и асимметрией вещества и антивещества во Вселенной. Тем самым электоральная статистика впервые включается в единую физическую картину мира, что применимо ко всем выборам, проведённым человечеством, и, согласно Теореме~3, ко всем, которые ещё будут проведены.
\end{zhplsignificance}

\begin{abstract}
\noindent Здесь мы впервые показываем, что положительная корреляция между явкой и долей голосов за партию «Единая Россия» на избирательных участках есть не социологическое, а гравитационное явление: захват участков компактным объектом, расположенным в точке «явка $100\,\%$, доля $100\,\%$». Проблема состоит в том, что двумерное распределение участков имеет форму кометы, хвост которой не описывается никаким равновесным распределением ядра и сопровождается избытком участков с круглыми значениями явки. Метод: вводится электоральное фазовое пространство с плоскими координатами $(u,o)$~--- долями списка, проголосовавшими за партию и за прочих; движение участка в поле аттрактора~$\At$ раскладывается на каналы аккреции и конверсии; спектр явки захваченного участка находится квантованием сопряжённой ей циклической переменной. Результаты: доказаны закон сохранения несогласных, теорема о прицельном параметре, электоральный принцип эквивалентности, теорема о квантовании явки и выполнение всех трёх условий Сахарова для электорогенеза. Показано, что явка захваченного участка определена по модулю $100\,\%$, что снимает парадокс явки $146\,\%$. Модельная реализация воспроизводит ядро, хвост и десятичную сверхструктуру спектра. Ограничения: согласно Теореме~3, все предсказания модели в точности совпадают с предсказаниями альтернативной гипотезы, не содержащей аттрактора.

\medskip
\noindent\textbf{Ключевые слова:} электоральная комета, компактный объект, аккреция, теорема Нётер, принцип эквивалентности, квантование явки, зона Бриллюэна, умклапп-процесс, условия Сахарова, лженаука
\end{abstract}

\section{Введение}

Каждый избирательный участок по итогам голосования характеризуется двумя числами: явкой $t$~--- долей избирателей списка, получивших бюллетени,~--- и долей $x$ голосов, поданных за партию «Единая Россия» (далее~--- партия~Е). Если нанести все участки страны на плоскость $(t,x)$, возникает устойчивая картина, известная в литературе как \emph{электоральная комета}~\cite{akkr}: компактное ядро в области умеренных значений $t$ и $x$ и протяжённый хвост, уходящий из ядра в правый верхний угол, к точке $(100\,\%,\,100\,\%)$. Внутри хвоста явка и доля партии~Е положительно коррелированы. Сверх того, распределение участков по явке в хвосте содержит резкие пики на целых значениях процента, особенно выраженные на значениях, кратных пяти~\cite{krug}.

Бытовое объяснение этих фактов исчерпывается рядом выражений, которые мы, следуя практике журнала, не приводим. В литературе известно также аддитивное объяснение~\cite{urn}, которое будет рассмотрено в разделе~4. Все они систематически игнорируют базовое наблюдение: хвост кометы направлен в одну точку. Траектории, сходящиеся в одну точку, в физике имеют стандартную интерпретацию~--- падение на компактный объект. В настоящей работе мы последовательно развиваем эту интерпретацию.

Статья построена следующим образом. В разделе~2 вводится электоральное фазовое пространство и описывается ядро кометы. В разделе~3 вводится аттрактор~$\At$ и доказываются закон сохранения несогласных и теорема о прицельном параметре. Раздел~4 посвящён электоральному принципу эквивалентности, раздел~5~--- квантованию явки, раздел~6~--- электорогенезу. Ограничения обсуждаются в разделе~7.

\section{Электоральное фазовое пространство}

\subsection{Плоские координаты}

\begin{definition}[Электоральная плоскость]
Состоянием участка с $N$ избирателями в списке называется пара $(u,o)$, где $u = U/N$~--- доля списка, проголосовавшая за партию~Е, $o = O/N$~--- доля списка, проголосовавшая за прочие партии. Физической областью является \emph{электоральный симплекс} $u\ge0$, $o\ge0$, $u+o\le1$. Наблюдаемые явка и доля партии выражаются через состояние как
\begin{equation}
t = u + o, \qquad x = \frac{u}{u+o}.
\label{eq:tx}
\end{equation}
\end{definition}

Координаты $(u,o)$ мы называем плоскими: в них, как будет показано ниже, движение участков прямолинейно. Координаты $(t,x)$, в которых традиционно строится комета, криволинейны.

\begin{lemma}[О координатной сингулярности нулевой явки]
Якобиан перехода $(t,x)\mapsto(u,o)$ равен $\bigl|\partial(u,o)/\partial(t,x)\bigr| = t$ и обращается в нуль при $t=0$, где доля $x$ не определена. Эта сингулярность является координатной и устраняется переходом к плоским координатам.
\end{lemma}

\begin{proof}
Из~\eqref{eq:tx} $u = tx$, $o = t(1-x)$, откуда $\det\begin{pmatrix} x & t\\ 1-x & -t\end{pmatrix} = -t$. В координатах $(u,o)$ точка $t=0$ есть регулярная вершина симплекса.
\end{proof}

Лемма~1 является электоральным аналогом сингулярности полярных координат в начале отсчёта: участок с нулевой явкой физически существует, но не имеет определённой политической ориентации.

\subsection{Ядро кометы}

\begin{definition}[Свободный участок]
Участок называется \emph{свободным}, если его состояние распределено по закону Гиббса с гауссовой плотностью в плоскости $(t,x)$ с центром $(t_0,x_0)$ и дисперсиями $\sigma_t^2$, $\sigma_x^2$. Множество свободных участков образует \emph{ядро} кометы.
\end{definition}

Гауссова форма ядра не требует специальных предположений: она следует из центральной предельной теоремы для суммы независимых решений избирателей и была бы наблюдаема при любых выборах, если бы все участки были свободными. В дальнейшем принимаются характерные значения $t_0 = x_0 = 0{,}40$, $\sigma_t = 0{,}07$, $\sigma_x = 0{,}06$; в плоских координатах центру ядра соответствует $u_0 = 0{,}16$, $o_0 = 0{,}24$.

Модельная реализация кометы при этих параметрах показана на рис.~\ref{fig:comet}. Несвободные участки, образующие хвост, будут описаны в разделе~3.

\begin{figure}[!htbp]
\centering
\begin{tikzpicture}
\begin{axis}[
  width=0.72\linewidth, height=0.64\linewidth,
  xlabel={явка $t$, \%}, ylabel={доля партии Е $x$, \%},
  xmin=0, xmax=100, ymin=0, ymax=100,
  xtick={0,20,...,100}, ytick={0,20,...,100},
  grid=major,
  legend style={at={(0.03,0.97)}, anchor=north west, font=\footnotesize},
  tick label style={font=\footnotesize}, label style={font=\footnotesize}]
  \addplot[only marks, mark=*, mark size=0.55pt, blue!70!black, opacity=0.55] table[x=t,y=x] {\tabFree};
  \addlegendentry{свободные участки (ядро)}
  \addplot[only marks, mark=*, mark size=0.55pt, red!75!black, opacity=0.7] table[x=t,y=x] {\tabCap};
  \addlegendentry{захваченные участки (хвост)}
  \addplot[only marks, mark=star, mark size=4pt, very thick, green!45!black] coordinates {(100,100)};
  \addlegendentry{аттрактор $\At$}
\end{axis}
\end{tikzpicture}
\caption{Электоральная комета: модельная реализация для $2600$ участков (параметры ядра~--- раздел~2.2, параметры захвата~--- раздел~3, квантования~--- раздел~5). Каждая точка~--- участок; число избирателей в списке равномерно распределено от $800$ до $2499$, голоса целочисленны. Хвост направлен в точку $\At$; вертикальные «столбики» в хвосте при круглых значениях явки объясняются в разделе~5.}
\label{fig:comet}
\end{figure}

\section{Аттрактор и закон сохранения несогласных}

\subsection{Каналы захвата}

\begin{definition}[Электоральный аттрактор]
Электоральным аттрактором называется компактный объект $\At$, расположенный в вершине симплекса $(u,o) = (1,0)$, что соответствует $(t,x) = (100\,\%,\,100\,\%)$. Участок, попавший в поле аттрактора (\emph{захваченный} участок), движется по закону
\begin{equation}
\frac{du}{d\lambda} = \alpha + \beta, \qquad \frac{do}{d\lambda} = -\beta, \qquad \alpha,\beta\ge0,
\label{eq:flow}
\end{equation}
где $\lambda$~--- аффинный параметр захвата. Слагаемое с $\alpha$ называется каналом \emph{аккреции} (голоса партии~Е возрастают, голоса прочих не изменяются), слагаемое с $\beta$~--- каналом \emph{конверсии} (голоса прочих переходят к партии~Е с сохранением явки). Отношение $\kappa = \beta/(\alpha+\beta)\in[0,1]$ называется \emph{коэффициентом конверсии}.
\end{definition}

Из~\eqref{eq:flow} видно, что в плоских координатах захваченный участок движется по прямой с наклоном $do/du = -\kappa$. Именно поэтому координаты $(u,o)$ и названы плоскими: геодезические в них~--- прямые. Криволинейность хвоста в координатах $(t,x)$ есть чисто координатный эффект.

\subsection{Чистая аккреция}

\begin{theorem}[Закон сохранения несогласных]
Пусть захват происходит по каналу чистой аккреции, т.\,е.\ движение участка описывается лагранжианом
\[
L = \tfrac12\bigl(\dot u^2 + \dot o^2\bigr) - \Phi(u),
\]
где потенциал аттрактора $\Phi$ зависит только от $u$, и до захвата участок покоился ($\dot o(0) = 0$). Тогда доля несогласных $b = o = t(1-x)$ является интегралом движения, а траектория участка в плоскости $(t,x)$ есть гипербола
\begin{equation}
x(t) = 1 - \frac{b}{t}.
\label{eq:noether}
\end{equation}
\end{theorem}

\begin{proof}
Лагранжиан инвариантен относительно сдвигов $o\mapsto o+\varepsilon$, поэтому по теореме Нётер сохраняется сопряжённый импульс $p_o = \dot o$. Из начального условия $\dot o\equiv0$, т.\,е.\ $o\equiv o(0)=b$. Подставляя $o=b$ в~\eqref{eq:tx}, получаем $x = (t-b)/t$.
\end{proof}

Физический смысл Теоремы~1 состоит в том, что при аккреции абсолютное число голосов за прочие партии не зависит от явки: весь прирост явки в хвосте целиком приходится на партию~Е. Гиперболы~\eqref{eq:noether} мы называем \emph{гиперболами Нётер} (рис.~\ref{fig:flow}, штриховые линии).

\begin{theorem}[О прицельном параметре]
Траектория чистой аккреции~\eqref{eq:noether} достигает границы симплекса $t=1$ в точке $x_{\max} = 1-b$ и проходит мимо аттрактора на расстоянии $b$, которое мы называем \emph{прицельным параметром}. Аттрактор достижим тогда и только тогда, когда $b=0$. Участок, стартовавший из точки ядра $(u_0,o_0)$ с $o_0>0$, попадает в $\At$ в том и только в том случае, когда его коэффициент конверсии равен
\begin{equation}
\kappa^* = \frac{o_0}{1-u_0}.
\label{eq:kstar}
\end{equation}
\end{theorem}

\begin{proof}
Первое утверждение следует из~\eqref{eq:noether} при $t=1$. Для второго: прямая $o = o_0 - \kappa(u-u_0)$ проходит через $(1,0)$ тогда и только тогда, когда $0 = o_0 - \kappa(1-u_0)$.
\end{proof}

\begin{corollary}
Для центра ядра $u_0 = 0{,}16$, $o_0 = 0{,}24$ имеем $\kappa^* = 2/7$. Следовательно, из каждых семи голосов, приобретённых партией~Е при радиальном падении участка на аттрактор, пять поступают по каналу аккреции, а два~--- по каналу конверсии. Наблюдение участков в окрестности точки $(100\,\%,\,100\,\%)$ доказывает существование канала конверсии, т.\,е.\ процесса, при котором голоса прочих партий переходят к партии~Е после того, как были поданы.
\end{corollary}

Траектории радиального падения в координатах $(t,x)$ получаются исключением $\lambda$ из~\eqref{eq:flow} при $\kappa=\kappa^*$:
\begin{equation}
x(t) = 1 - \frac{o_0}{1-t_0}\cdot\frac{1-t}{t},
\label{eq:radial}
\end{equation}
и все они, независимо от точки старта, сходятся в $\At$ (рис.~\ref{fig:flow}, сплошные линии). Реальный хвост кометы заполнен траекториями с промежуточными $0\le\kappa\le\kappa^*$; именно такой захват использован в модельной реализации рис.~\ref{fig:comet}.

\begin{figure}[!htbp]
\centering
\begin{tikzpicture}
\begin{axis}[
  width=0.72\linewidth, height=0.64\linewidth,
  xlabel={явка $t$, \%}, ylabel={доля партии Е $x$, \%},
  xmin=0, xmax=100, ymin=0, ymax=100,
  xtick={0,20,...,100}, ytick={0,20,...,100},
  grid=major,
  legend style={at={(0.97,0.03)}, anchor=south east, font=\footnotesize},
  legend cell align=left, clip=false,
  tick label style={font=\footnotesize}, label style={font=\footnotesize}]
  % ядро: эллипсы 1 sigma и 2 sigma
  \addplot[thick, blue!70!black, fill=blue!70!black, fill opacity=0.12, domain=0:360, samples=90]
    ({40+7*cos(x)},{40+6*sin(x)});
  \addlegendentry{ядро, $1\sigma$ и $2\sigma$}
  \addplot[thick, blue!70!black, dotted, domain=0:360, samples=90, forget plot]
    ({40+14*cos(x)},{40+12*sin(x)});
  % гиперболы Нётер x = 1 - b/t, старт из точек ядра
  \addplot[thick, dashed, orange!85!black, domain=40:100, samples=80] {100*(1-0.24/(x/100))};
  \addlegendentry{гиперболы Нётер, $\kappa=0$}
  \addplot[thick, dashed, orange!85!black, domain=33:100, samples=80, forget plot] {100*(1-0.16/(x/100))};
  \addplot[thick, dashed, orange!85!black, domain=47:100, samples=80, forget plot] {100*(1-0.329/(x/100))};
  % радиальное падение x = 1 - c(1-t)/t, c = o0/(1-t0)
  \addplot[very thick, red!75!black, domain=40:100, samples=80] {100*(1-0.4*(1-x/100)/(x/100))};
  \addlegendentry{радиальное падение, $\kappa=\kappa^*$}
  \addplot[very thick, red!75!black, domain=33:100, samples=80, forget plot] {100*(1-(0.16/0.67)*(1-x/100)/(x/100))};
  \addplot[very thick, red!75!black, domain=47:100, samples=80, forget plot] {100*(1-(0.329/0.53)*(1-x/100)/(x/100))};
  % точки старта
  \addplot[only marks, mark=*, mark size=1.8pt, black, forget plot] coordinates {(40,40) (33,51.5) (47,30)};
  % концы гипербол
  \addplot[only marks, mark=o, mark size=2pt, orange!85!black, forget plot] coordinates {(100,76) (100,84) (100,67.1)};
  \addplot[only marks, mark=star, mark size=4pt, very thick, green!45!black] coordinates {(100,100)};
  \addlegendentry{аттрактор $\At$}
  \draw[<->, thick, xshift=5pt] (axis cs:100,76) -- (axis cs:100,100) node[midway, right, font=\footnotesize] {$b$};
\end{axis}
\end{tikzpicture}
\caption{Траектории захвата для трёх точек старта в ядре ($\bullet$): $(t,x) = (40, 40)$, $(33; 51{,}5)$, $(47, 30)\,\%$, т.\,е.\ $b = o_0 = 0{,}24$; $0{,}16$; $0{,}329$. Штриховые линии~--- гиперболы Нётер~\eqref{eq:noether} (чистая аккреция), оканчивающиеся на $t=100\,\%$ в точках $x = 1-b$ ($\circ$); стрелкой показан прицельный параметр центральной траектории. Сплошные линии~--- радиальное падение~\eqref{eq:radial}, сходящееся в аттракторе.}
\label{fig:flow}
\end{figure}

\section{Электоральный принцип эквивалентности}

В литературе предложено объяснение хвоста кометы, не требующее аттрактора~\cite{urn}. Согласно этой \emph{аддитивной гипотезе}, на участке в урну дополнительно помещаются бюллетени, заполненные за партию~Е, в количестве $\alpha N$, а $\beta N$ уже поданных бюллетеней за прочие партии учитываются в пользу партии~Е. Покажем, что эта гипотеза не может быть отличена от гравитационной.

\begin{theorem}[Электоральный принцип эквивалентности]
Пусть $\mathcal{P}$~--- протокол участка, т.\,е.\ любая функция от $(N, U, O)$. Тогда для любых $\alpha,\beta\ge0$ протокол участка, захваченного аттрактором $\At$ по закону~\eqref{eq:flow} на единичном отрезке аффинного параметра, совпадает с протоколом того же участка при аддитивной гипотезе с теми же $\alpha$, $\beta$. Никакое измерение, выполненное над протоколом, не позволяет отличить гравитационный захват от добавления бюллетеней в урну.
\end{theorem}

\begin{proof}
Интегрирование~\eqref{eq:flow} по $\lambda\in[0,1]$ даёт отображение $(u,o)\mapsto(u+\alpha+\beta,\ o-\beta)$. Аддитивная гипотеза по определению даёт $(U,O)\mapsto(U+\alpha N+\beta N,\ O-\beta N)$, т.\,е.\ то же отображение. Протокол есть функция от $(N,U,O)$ и потому совпадает.
\end{proof}

Теорема~3 является точным электоральным аналогом принципа эквивалентности Эйнштейна: наблюдатель в закрытом лифте не может отличить тяготение от ускорения, наблюдатель с протоколом~--- тяготение аттрактора от добавления бюллетеней (рис.~\ref{fig:urn}).

\begin{figure}[!htbp]
\centering
\begin{tikzpicture}[scale=0.9,
  bal/.style={draw, fill=\zhpldark{gray!10}, minimum width=0.55cm, minimum height=0.18cm, inner sep=0pt},
  balE/.style={draw, fill=\zhpldark{red!25}, minimum width=0.55cm, minimum height=0.18cm, inner sep=0pt},
  every node/.append style={font=\footnotesize}]
  \foreach \cx in {0,6.2} {
    \draw[very thick] (\cx-0.45,1.1) -- (\cx-1.3,1.1) -- (\cx-1.3,-1.1) -- (\cx+1.3,-1.1) -- (\cx+1.3,1.1) -- (\cx+0.45,1.1);
    \foreach \y in {-0.9,-0.65,-0.4} \node[bal] at (\cx-0.6,\y) {};
    \foreach \y in {-0.9,-0.65,-0.4,-0.15,0.1} \node[balE] at (\cx+0.6,\y) {};
  }
  % левая урна: поле аттрактора
  \draw[-{Latex[length=3mm]}, very thick, green!45!black] (1.6,-0.3) -- (2.6,0.7) node[right] {к $\At$};
  \node[align=center] at (0,-1.75) {гравитационный захват\\ по закону~\eqref{eq:flow}};
  % правая урна: аддитивная гипотеза
  \node[balE, rotate=90] at (6.1,1.5) {};
  \node[balE, rotate=90] at (6.3,2.25) {};
  \node[anchor=west] at (6.6,1.85) {$+\alpha N$};
  \node[align=center] at (6.2,-1.75) {аддитивная гипотеза\\ с теми же $\alpha$, $\beta$};
  % общий протокол
  \node[draw, rounded corners, fill=\zhpldark{green!8}, align=center] (P) at (3.1,-3.3)
    {протокол: $U+(\alpha+\beta)N,\quad O-\beta N$};
  \draw[-{Latex[length=2.4mm]}, thick] (0,-2.3) -- (P);
  \draw[-{Latex[length=2.4mm]}, thick] (6.2,-2.3) -- (P);
\end{tikzpicture}
\caption{Урна Эйнштейна. Слева~--- участок в поле аттрактора $\At$, справа~--- тот же участок при аддитивной гипотезе. Бюллетени за партию~Е показаны цветом. Оба случая дают один и тот же протокол (Теорема~3), и никакой эксперимент внутри участка не позволяет установить, в каком из них находится наблюдатель.}
\label{fig:urn}
\end{figure}

Как и в общей теории относительности, эквивалентность локальна. В масштабе всей страны гравитационная гипотеза предсказывает приливные эффекты~--- зависимость темпа захвата от расстояния участка до аттрактора,~--- тогда как аддитивная гипотеза их не предсказывает. К сожалению, проверка требует размещения наблюдателей в непосредственной окрестности точки $(100\,\%,\,100\,\%)$, куда наблюдатели, как правило, не допускаются.

\section{Квантование явки}

\subsection{Связанные состояния}

Захваченный участок, в отличие от свободного, совершает финитное движение в поле аттрактора. Для финитного движения существуют переменные действие--угол; обозначим через $\varphi\in[0,2\pi)$ циклическую переменную захваченного участка, которую мы называем \emph{фазой протокола}. Явка является сопряжённым к ней импульсом и в квантовой теории представляется оператором
\begin{equation}
\hat T = -i\hb\,\frac{\partial}{\partial\varphi}, \qquad \psi(\varphi+2\pi) = \psi(\varphi),
\label{eq:T}
\end{equation}
где $\hb$~--- \emph{электоральная постоянная Планка}.

\begin{theorem}[О квантовании явки]
Спектр оператора~\eqref{eq:T} на захваченном участке дискретен:
\[
t_n = n\hb, \qquad n = 0, 1, \dots, 100, \qquad \hb = 1\,\%,
\]
с собственными функциями $\psi_n = e^{in\varphi}/\sqrt{2\pi}$. На свободном участке циклическая переменная отсутствует, и спектр явки непрерывен.
\end{theorem}

\begin{proof}
Из~\eqref{eq:T} $\hat T\psi_n = n\hb\,\psi_n$, а периодическое граничное условие требует $e^{2\pi i n} = 1$, т.\,е.\ $n\in\mathbb{Z}$. Ограничение $0\le n\le100$ следует из $0\le t\le1$. Для свободного участка движение инфинитно, граничное условие отсутствует, и спектр совпадает со спектром оператора импульса на прямой.
\end{proof}

Значение $\hb = 1\,\%$ установлено по положению пиков (рис.~\ref{fig:spectrum}). Теорема~4 объясняет центральный наблюдательный факт~\cite{krug}: пики на целых значениях явки наблюдаются только в хвосте кометы, т.\,е.\ у захваченных участков, тогда как распределение ядра гладко. Состояния с $n>100$ лежат, на первый взгляд, вне физической области; они рассматриваются в разделе~5.3.

\subsection{Десятичная сверхструктура}

Уровни, кратные пяти, заселены заметно сильнее прочих, а кратные десяти~--- сильнее кратных пяти. Для описания этой \emph{десятичной сверхструктуры} вводится \emph{гамильтониан округления}
\begin{equation}
\hat H_{\text{окр}} = -J_5\,\hat\Pi_5 - J_{10}\,\hat\Pi_{10},
\label{eq:Hr}
\end{equation}
где $\hat\Pi_k$~--- проектор на подпространство состояний $\psi_n$ с $n$, кратным $k$, а $J_5, J_{10}>0$~--- энергии связи круглого числа. Поскольку $[\hat H_{\text{окр}}, \hat T] = 0$, собственные состояния явки являются и собственными состояниями~\eqref{eq:Hr}, и в равновесии при электоральной температуре $\Theta$ заселённости уровней подчиняются распределению Гиббса:
\begin{equation}
p_n \propto g_n = \exp\!\Bigl[\tfrac{J_5}{\Theta}\,[5\,|\,n] + \tfrac{J_{10}}{\Theta}\,[10\,|\,n]\Bigr],
\end{equation}
где $[k\,|\,n]$ равно единице при $n$, кратном $k$, и нулю в противном случае. Значения $J_5/\Theta = \ln4$, $J_{10}/\Theta = \ln2$ дают отношение заселённостей
\begin{equation}
g_{\text{прочие}} : g_{5} : g_{10} = 1 : 4 : 8,
\end{equation}
которое и использовано в модельной реализации (рис.~\ref{fig:spectrum}).

\begin{figure}[!htbp]
\centering
\begin{tikzpicture}
\begin{axis}[
  width=0.95\linewidth, height=5.8cm,
  xlabel={явка $t$, \%}, ylabel={число участков в интервале $0{,}25\,\%$},
  xmin=20, xmax=100.5, ymin=0, ymax=780,
  xtick={20,30,...,100}, minor xtick={25,35,...,95},
  grid=major,
  legend style={at={(0.375,0.97)}, anchor=north west, font=\footnotesize},
  legend cell align=left,
  tick label style={font=\footnotesize}, label style={font=\footnotesize}]
  \addplot[const plot, draw=none, fill=red!75!black, fill opacity=0.35, area legend] table[x=t,y=cap] {\tabHist} \closedcycle;
  \addlegendentry{захваченные участки}
  \addplot[const plot, thin, blue!70!black] table[x=t,y=all] {\tabHist};
  \addlegendentry{все участки}
\end{axis}
\end{tikzpicture}
\caption{Спектр явки: модельная гистограмма для $60\,000$ участков, ширина интервала $0{,}25\,\%$, центры интервалов кратны $0{,}25\,\%$. Широкий максимум около $40\,\%$~--- непрерывный спектр свободных участков ядра. Узкие линии~--- дискретный спектр захваченных участков (Теорема~4); линии при $n$, кратном $10$, и при $n$, кратном $5$, усилены гамильтонианом округления~\eqref{eq:Hr}. Линии при остальных целых $n$ слабее, но также различимы.}
\label{fig:spectrum}
\end{figure}

Отметим одно важное следствие. Доля партии $x$ также может быть сопоставлена сопряжённой циклической переменной второй фазы протокола; соответствующий оператор $\hat X$ коммутирует с $\hat T$, так как действует на другую переменную. Поэтому существуют совместные собственные состояния, в которых \emph{одновременно} круглыми являются и явка, и доля партии~Е. Их наблюдение~\cite{krug} служит прямым подтверждением того, что фаза протокола двумерна.

Наконец, десятичность сверхструктуры не является случайной. При пересчёте наблюдаемых явок в двенадцатеричную систему счисления выделенные уровни не переходят в круглые двенадцатеричные значения. Следовательно, квантование происходит именно в десятичной системе~--- единственной, используемой на участках.

\subsection{Явка $146\,\%$ и зона Бриллюэна}

Ограничение $n\le100$ в Теореме~4 долгое время считалось очевидным. Однако в эфире одного из федеральных телеканалов в 2011~году было зарегистрировано состояние с $t = 146\,\%$, которое не может быть собственным значением оператора~\eqref{eq:T} на отрезке $0\le t\le1$. Покажем, что противоречие снимается учётом решёточной структуры фазы протокола.

В связанном состоянии протокол фиксирует явку в целых процентах. Следовательно, фаза протокола принимает не континуум, а лишь $100$ значений $\varphi_k = 2\pi k/100$, $k = 0,\dots,99$, образуя одномерную решётку с периодом $a = 2\pi/100$. По теореме Блоха импульс, сопряжённый решёточной координате, определён лишь с точностью до вектора обратной решётки:
\begin{equation}
t \sim t + G, \qquad G = \frac{2\pi\hb}{a} = 100\,\hb = 100\,\%.
\label{eq:umklapp}
\end{equation}

\begin{corollary}[Об умклапп-процессе]
Явка захваченного участка определена лишь по модулю $100\,\%$, а физически различимые значения образуют первую зону Бриллюэна $0\le t<100\,\%$. В частности, состояние $t = 146\,\%$ физически эквивалентно состоянию $t = 46\,\%$.
\end{corollary}

\begin{proof}
Плоские волны $e^{i t\varphi/\hb}$ и $e^{i(t+G)\varphi/\hb}$ совпадают во всех узлах решётки, так как $e^{iG\varphi_k/\hb} = e^{2\pi i k} = 1$. Следовательно, никакое измерение на решётке их не различает. Приведение в первую зону даёт $146 - 100 = 46$.
\end{proof}

Приведённое значение $46\,\%$ отстоит от центра ядра $t_0 = 40\,\%$ менее чем на $1\sigma_t$ (рис.~\ref{fig:bz}). Таким образом, наблюдавшееся в эфире состояние после приведения в первую зону оказывается совершенно типичным \emph{свободным} участком. Переход $146\,\%\to46\,\%$ мы называем электоральным умклапп-процессом; по-видимому, именно он происходит при передаче данных из эфира в итоговый протокол. Распространённая интерпретация, согласно которой $146\,\%$ было не явкой, а суммой долей нескольких партий, для настоящей теории несущественна: доли партий также являются решёточными импульсами и также определены по модулю $100\,\%$.

\begin{figure}[!htbp]
\centering
\begin{tikzpicture}[x=0.066cm, y=1cm, every node/.append style={font=\footnotesize}]
  \fill[\zhpldark{green!10}] (0,-0.28) rectangle (100,0.28);
  \fill[blue!70!black, opacity=0.25] (33,-0.18) rectangle (47,0.18);
  \draw[-{Latex[length=2mm]}, thick] (-4,0) -- (206,0) node[right] {$t$, \%};
  \foreach \v in {0,20,...,200} {\draw (\v,-0.08) -- (\v,0.08); \node[anchor=north, font=\scriptsize] at (\v,-0.1) {\v};}
  \draw[dashed] (100,-0.45) -- (100,0.45);
  \node[blue!70!black, anchor=north east] at (48,-0.45) {ядро, $t_0\pm\sigma_t$};
  \node[green!40!black, anchor=north west] at (53,-0.45) {первая зона Бриллюэна};
  \fill[red!75!black] (146,0) circle (2.4pt);
  \node[red!75!black, anchor=north] at (146,-0.45) {$146\,\%$ (эфир)};
  \fill[red!75!black] (46,0) circle (2.4pt);
  \draw[-{Latex[length=2.4mm]}, very thick, red!75!black] (146,0.12) to[bend right=28] node[above] {$-G = -100\,\%$} (46,0.12);
  \draw[-{Latex[length=2mm]}, thick, densely dashed, green!45!black] (100,-1.0) to[bend left=8] node[below] {$100\,\%\sim0\,\%$ (Следствие~3)} (0,-1.0);
\end{tikzpicture}
\caption{Приведение явки в первую зону Бриллюэна по формуле~\eqref{eq:umklapp}. Состояние $146\,\%$ переводится вектором обратной решётки $-G$ в $46\,\%$, лежащее в пределах $1\sigma_t$ от центра ядра кометы. Штриховая стрелка~--- отождествление $100\,\%\sim0\,\%$.}
\label{fig:bz}
\end{figure}

\begin{corollary}[О стопроцентной явке]
Участок с явкой $100\,\%$ физически эквивалентен участку, на который не пришёл ни один избиратель. В частности, аттрактор $\At$ при приведении в первую зону совпадает с координатной сингулярностью нулевой явки (Лемма~1): это один и тот же объект, наблюдаемый из соседних зон Бриллюэна.
\end{corollary}

\begin{proof}
Формула~\eqref{eq:umklapp} при $t = 100\,\%$ даёт $t\sim0$. Доля $x$ при $t=0$ не определена (Лемма~1), что согласуется с тем, что у аттрактора она равна $100\,\%$ лишь в одной из зон.
\end{proof}

Следствие~3 является единственным предсказанием настоящей работы, допускающим прямую экспериментальную проверку: для этого достаточно опросить избирателей участка со стопроцентной явкой.

\section{Электорогенез}

Ядро кометы симметрично относительно замены партии~Е на прочие партии (зарядового сопряжения $C: x\mapsto1-x$) в том смысле, что гауссова форма ядра не выделяет ни одной из них. Хвост кометы эту симметрию нарушает: аттрактор расположен при $x=1$, тогда как сопряжённая ему точка $(t,x)=(100\,\%,\,0\,\%)$ участков не притягивает. Динамическое возникновение такой асимметрии в физике известно как бариогенез; мы называем его \emph{электорогенезом}.

\begin{theorem}[Об условиях Сахарова для электорогенеза]
В модели разделов 3--5 выполнены все три необходимых условия динамической генерации асимметрии между партией~Е и её зарядово-сопряжённым партнёром (табл.~\ref{tab:sakharov}).
\end{theorem}

\begin{proof}
(i) Канал аккреции увеличивает $U$ при неизменном $O$ и тем самым меняет полное число поданных голосов без соответствующего изменения числа избирателей, пришедших на участок в свободном состоянии. (ii) Отображение~\eqref{eq:flow} при $\alpha+\beta>0$ не коммутирует с $C$, а аттрактор не имеет сопряжённого партнёра. (iii) Распределение участков в хвосте не является распределением Гиббса ядра: хвост не имеет максимума и заканчивается в граничной точке симплекса.
\end{proof}

\begin{table}[!htbp]
\centering
\small
\begin{tabular}{@{}p{4.6cm}p{6.2cm}c@{}}
\toprule
Условие Сахарова & Электоральный аналог & Выполнено \\
\midrule
Несохранение барионного числа & Несохранение числа бюллетеней (канал аккреции) & да \\
Нарушение $C$ и $CP$ & Аттрактор при $x=1$ без партнёра при $x=0$ & да \\
Отклонение от теплового равновесия & Хвост не описывается гиббсовым распределением ядра & да \\
\bottomrule
\end{tabular}
\caption{Условия Сахарова в электоральной динамике.}
\label{tab:sakharov}
\end{table}

Теорема~5 показывает, что наблюдаемое преобладание партии~Е в хвосте кометы не требует начальной асимметрии: оно генерируется динамически, так же как, согласно современным представлениям, в ранней Вселенной генерировалось преобладание вещества над антивеществом.

\section{Обсуждение и ограничения}

Развитая теория описывает все три наблюдаемых свойства электоральной кометы~--- ядро, направленный хвост и десятичную структуру спектра явки~--- исходя из единственного допущения о существовании аттрактора~$\At$. Тем не менее необходимо честно указать на следующие ограничения.

\emph{Бритва Оккама.} Согласно Теореме~3, аддитивная гипотеза~\cite{urn} даёт в точности те же предсказания для любых наблюдаемых величин. При этом она не содержит аттрактора, т.\,е.\ требует на одну сущность меньше. Поскольку из двух эмпирически эквивалентных теорий принято предпочитать более экономную, аддитивная гипотеза строго предпочтительнее гравитационной, а вся развитая в настоящей работе теория избыточна. Это можно рассматривать как единственный строгий результат работы.

\emph{Нестационарность аттрактора.} Аттрактор проявляет себя исключительно в день голосования, а в промежутках между выборами не наблюдается. Это несовместимо с законом сохранения энергии и требует либо испарения аттрактора за время порядка избирательного цикла, либо предположения, что он существует постоянно, но взаимодействует с участками лишь в определённые дни. Выбор между этими возможностями выходит за рамки настоящей работы.

\emph{Происхождение асимметрии.} Теорема~5 устанавливает, что условия электорогенеза выполнены, но не объясняет, почему аттрактор расположен именно при $x=1$ для партии~Е. Как и в случае барионной асимметрии Вселенной, этот вопрос остаётся открытым.

\section{Заключение}

Корреляция явки и доли голосов за партию «Единая Россия» на избирательных участках получила строгое физическое объяснение: захват участков компактным аттрактором с сохранением несогласных, квантование явки в связанных состояниях и динамический электорогенез. Все наблюдаемые следствия теории совпадают с наблюдаемыми следствиями альтернативного объяснения, что делает теорию одновременно полной и ненужной. Авторы считают это достаточным основанием для дальнейшего финансирования.

\section*{Вклад авторов}
Автор~1: концептуализация, методология, написание -- черновик.
Автор~2: программное обеспечение, визуализация, формальный анализ, написание -- рецензирование и редактирование, техническая поддержка бесконечного терпения.

\section*{Финансирование}
Исследование выполнено при поддержке гранта Института Прикладной Метафизики № 2026-146-$\infty$ «Теория великого единства».

\section*{Доступность данных}
Модельные данные и код их генерации доступны по обоснованному запросу -- а также по необоснованному, с той же вероятностью успеха. Протоколы участков, находящихся в связанных состояниях, доступны в виде целых процентов.

\begin{thebibliography}{9}
\bibitem{akkr} Аккрецов,~А.~А. \textit{О захвате избирательных участков компактными объектами.} Труды Института Прикладной Метафизики, б.г.
\bibitem{urn} Урнов,~У.~У. \textit{Аддитивная модель протокола: урна как открытая система.} Неопубликованная рукопись.
\bibitem{krug} Круглов,~К.~К. \textit{Десятичная сверхструктура спектра явки и совместные круглые состояния.} Препринт, представлен ровно в $12{:}00$.
\end{thebibliography}

\noindent\textbf{Конфликт интересов.} Не заявлен, хотя, согласно Теореме~3, никакое измерение внутри редакции не позволяет отличить его отсутствие от присутствия.

\medskip
\noindent\textbf{Примечание редакции.} Авторы пожелали остаться анонимными по причинам, не имеющим отношения к содержанию работы. Их явка на церемонию вручения премии Института составила $0\,\%$, что, согласно Следствию~3, эквивалентно $100\,\%$.

\end{document}
